% Consolidated 39-nuclear-reactions from ce5faa3e9e459c6fa8c4f0ed5dbc763b44d938dd; original TeX fragments preserved.

% === BEGIN 39-nuclear-reactions-249-258.tex blob 1cffb11cd68d61e7c04a19d525ede731d3c3ea50 ===
% 原书第六篇6.6节249--258，印刷260--261，原PDF269--270；
% 书末答案349--350，PDF358--359。
\nuclearreactionproblem{249}{260--261}
动能 $T_\alpha=7.0\,\mathrm{MeV}$ 的
$\alpha$ 粒子与起初静止的 $\mathrm{^{6}Li}$ 核
弹性散射。原题称“两粒子出射方向夹角”
$\theta=60^\circ$，求反冲锂核动能。
\exsolution
令入射粒子质量 $m=4u$、靶核质量 $M=6u$，
散射后两个粒子动量模为 $p'_a,p_R$，
它们之间的实际夹角为 $\gamma$。
动量、能量守恒分别为
\[
 p_0^2=p_a'^2+p_R^2+2p_a'p_R\cos\gamma,
 \qquad
 \frac{p_0^2}{2m}=
 \frac{p_a'^2}{2m}+\frac{p_R^2}{2M}.
\]
两式消去 $p_0^2$：
\[
 2p_a'p_R\cos\gamma=
 \left(\frac mM-1\right)p_R^2.
\]
因 $m<M$，实际 $\cos\gamma<0$，
所以两出射动量夹角必须\emph{大于}$90^\circ$。
若原题角度 $60^\circ$ 是实际出射夹角，
则不存在物理解。
原书答案 $T_R=6.0\,\mathrm{MeV}$ 可由补角
$\gamma=120^\circ$ 求出：
利用
$p_a'^2=2m(T_\alpha-T_R)$、
$p_R^2=2MT_R$，得
\[
 \boxed{T_R=
 \frac{T_\alpha}{1+
 \frac{(M-m)^2}{4mM\cos^2\gamma}}=
 6.0\,\mathrm{MeV}}.
\]
\excheck \emph{印刷60°与解6.0MeV不同时满足带符号动量守恒}，
因为6MeV实际要求夹角120°。
本稿同时列出条件及书末结果，不悄悄修改出射角。

\nuclearreactionproblem{250}{261}
中子与起初静止的氘核弹性碰撞。
求中子损失动能占初始动能的比例：
(a)正面碰撞；(b)中子偏转 $90^\circ$。
\exsolution
令中子质量 $m$、氘核质量 $M\simeq2m$。
(a) 质量为$m$的粒子撞静止质量$M$物体，
弹性正碰后动能比例
$T'_n/T_n=[(m-M)/(m+M)]^2$，
所以
\[
 \boxed{\eta_a=
 1-\left(\frac{m-M}{m+M}\right)^2
 =\frac{4mM}{(m+M)^2}=\frac89}.
\]
(b) 末态中子与入射方向正交，
矢量动量守恒给
$p_R^2=p_0^2+p_n'^2$。
代能量守恒
$p_0^2/(2m)=p_n'^2/(2m)+p_R^2/(2M)$，
得到
$p_n'^2/p_0^2=(M-m)/(M+m)=1/3$，故
\[
 \boxed{\eta_b=1-\frac13=\frac23}.
\]
\excheck (b) $90^\circ$ 指中子偏转角，
不是中子和氘核两个出射粒子的夹角。

\nuclearreactionproblem{251}{261}
氘核与起初静止的质子弹性碰撞，
求氘核在实验室系中可能达到的最大散射角。
\exsolution
入射氘核质量$m_d\simeq2m_p$，
靶质子质量$M=m_p$。
质心系中入射氘核速度
$u^*=Mv_0/(m_d+M)$，
质心平移速度
$V=m_dv_0/(m_d+M)$。
实验室末态氘核速度矢量是
以 $V$ 为圆心、半径 $u^*$ 的速度圆。
从原点向该圆作切线时实验室偏转角最大：
\[
 \boxed{\sin\theta_{\max}=\frac{u^*}{V}
 =\frac M{m_d}\simeq\frac12},
 \quad
 \boxed{\theta_{\max}=30^\circ}.
\]
\excheck 重入射粒子撞轻靶核时，
其实验室偏角有限制，并非可取任意后散射角。

\nuclearreactionproblem{252}{261}
原书将核半径写为
$R=0.13A^{1/3}\,\mathrm{pm}$，
$A$ 为质量数。按此题面估算原子核质量密度
和核子数密度，并核对书末量级。
\exsolution
若严格采用原印
$r_0=0.13\,\mathrm{pm}=1.3\times10^{-13}\,\mathrm m$，
半径为$r_0A^{1/3}$，则
\[
 \rho=\frac{Au}{(4\pi/3)r_0^3A}
 =\frac{3u}{4\pi r_0^3},\quad
 n=\frac{3}{4\pi r_0^3}.
\]
数值约为
\[
 \boxed{\rho_{\rm literal}\simeq
 1.8\times10^{5}\,\mathrm{kg/cm^3}},\qquad
 \boxed{n_{\rm literal}\simeq
 1.1\times10^{32}\,\mathrm{nucleons/cm^3}}.
\]
这与书末给出的约
$2\times10^{11}\,\mathrm{kg/cm^3}$、
$10^{38}\,\mathrm{nucleons/cm^3}$
相差$10^6$。
若采用物理上公认的核尺度
$r_0=1.3\,\mathrm{fm}=0.0013\,\mathrm{pm}$，
则
\[
 \boxed{\rho_{\rm nuclear}\simeq
 2\times10^{11}\,\mathrm{kg/cm^3}},\quad
 \boxed{n_{\rm nuclear}\simeq
 10^{38}\,\mathrm{nucleons/cm^3}},
\]
恢复原书答案量级。
\excheck 原图确印$0.13A^{1/3}\,\mathrm{pm}$，
而从书末密度可判断物理上需要
$0.0013A^{1/3}\,\mathrm{pm}$；
本译稿\emph{并列两种数据}并列入勘误。

\nuclearreactionproblem{253}{261}
利用质量数和电荷数守恒补全：
(a) $\mathrm{^{10}B}(x,\alpha)\mathrm{^{8}Be}$；
(b) $\mathrm{^{17}O}(d,n)x$；
(c) $\mathrm{^{23}Na}(p,x)\mathrm{^{20}Ne}$；
(d) $x(p,n)\mathrm{^{37}Ar}$。
\exsolution
核反应的核子总数$A$和电荷总数$Z$守恒。
(a) 输入粒子应有$A=8+4-10=2$、
$Z=4+2-5=1$，即
$\boxed{x=d=\mathrm{^{2}H}}$。
(b) 输出
$A=17+2-1=18$、
$Z=8+1-0=9$，即
$\boxed{x=\mathrm{^{18}F}}$。
(c) 输入后$A=24,Z=12$，
$\mathrm{^{20}Ne}$有$A=20,Z=10$，
剩余 $\boxed{x=\alpha=\mathrm{^{4}He}}$。
(d) 输入靶核$A=37+1-1=37$、
$Z=18-1=17$，
故$\boxed{x=\mathrm{^{37}Cl}}$。
\excheck 原书简答(b)纯文本出现
$\mathrm{^{17}F}$，但题面明确是
$\mathrm{^{17}O}(d,n)$，
按质量数守恒产物必须是$\mathrm{^{18}F}$；
已记录底本可能的上标排印或识别差异。

\nuclearreactionproblem{254}{261}
证明核的结合能可按质量超额计算：
$E_b=[Z\Delta_H+(A-Z)\Delta_n-\Delta_{\rm at}]c^2$。
\exsolution
若用中性原子的实测质量$M_{\rm at}$，
核结合能（忽略电子化学结合能）为
\[
 E_b=[Zm_H+(A-Z)m_n-M_{\rm at}]c^2.
\]
把原子质量与中子质量写成
$m_H=u+\Delta_H$、
$m_n=u+\Delta_n$、
$M_{\rm at}=Au+\Delta_{\rm at}$。
则
\[
 Zm_H+(A-Z)m_n-M_{\rm at}
 =Z\Delta_H+(A-Z)\Delta_n-\Delta_{\rm at},
\]
其中$Zu+(A-Z)u-Au=0$，因此
\[
 \boxed{E_b=
 [Z\Delta_H+(A-Z)\Delta_n-\Delta_{\rm at}]c^2}.
\]
\excheck 选用\emph{氢原子}质量而非裸质子质量，
正是为了抵消相应的原子电子静质量。

\nuclearreactionproblem{255}{261}
某核的质子数与中子数相等，
且半径比 $\mathrm{^{27}Al}$ 核小1.5倍。
求核种与其结合能。
\exsolution
核半径比例$R\propto A^{1/3}$，
所以
\[
 \frac{R_x}{R_{\mathrm{Al}}}=\frac1{1.5}
 \Longrightarrow A_x=\frac{27}{1.5^3}=8.
\]
质子数等于中子数，故$Z=8/2=4$，
该核为$\boxed{\mathrm{^{8}Be}}$。
按原书附录质量超额
$\Delta_H=0.00783u$、
$\Delta_n=0.00867u$、
$\Delta_{\mathrm{^{8}Be}}=0.00531u$，
\[
 \boxed{E_b=
 [4(0.00783)+4(0.00867)-0.00531]
 (931.5\,\mathrm{MeV})
 \simeq56.5\,\mathrm{MeV}}.
\]
\excheck $\mathrm{^{8}Be}$ 对裂变为两$\alpha$
粒子是不稳定的，但仍可计算它的核结合能。

\nuclearreactionproblem{256}{261}
用原子质量表求：
(a) $\mathrm{^{16}O}$ 平均每核子结合能；
(b) $\mathrm{^{11}B}$ 中单中子及 $\alpha$
的分离能；(c) 把 $\mathrm{^{16}O}$ 分成四个
相同粒子所需能量。
\exsolution
原书质量超额（单位$u$）：
$\Delta_H=0.00783$、
$\Delta_n=0.00867$、
$\Delta_{\mathrm{^{16}O}}=-0.00509$、
$\Delta_{\mathrm{^{11}B}}=0.00930$、
$\Delta_{\mathrm{^{10}B}}=0.01294$、
$\Delta_{\mathrm{^{7}Li}}=0.01601$、
$\Delta_\alpha=0.00260$。
(a)
\[
 \boxed{\frac{E_b(\mathrm{^{16}O})}{16}
 =\frac{[8\Delta_H+8\Delta_n-\Delta_O]c^2}{16}
 \simeq8.0\,\mathrm{MeV}}.
\]
(b) 分离中子反应
$\mathrm{^{11}B}\to\mathrm{^{10}B}+n$，
以及分离$\alpha$反应
$\mathrm{^{11}B}\to\mathrm{^{7}Li}+\alpha$，故
\[
 \boxed{S_n=(\Delta_{\mathrm{^{10}B}}
 +\Delta_n-\Delta_{\mathrm{^{11}B}})c^2
 \simeq11.5\,\mathrm{MeV}},
\]
\[
 \boxed{S_\alpha=(\Delta_{\mathrm{^{7}Li}}
 +\Delta_\alpha-\Delta_{\mathrm{^{11}B}})c^2
 \simeq8.7\,\mathrm{MeV}}.
\]
(c) 四个相同粒子应是$\alpha$，
\[
 \boxed{E= [4\Delta_\alpha-\Delta_O]c^2
 \simeq14.4\,\mathrm{MeV}\approx14.5\,\mathrm{MeV}}.
\]
\excheck 分离能是末态总静能减母核静能；
守恒的核子总数使整数原子质量数自动抵消。

\nuclearreactionproblem{257}{261}
求 $\mathrm{^{11}B}$ 核中中子与质子的
分离能之差，并解释其原因。
\exsolution
中子分离能
$S_n=[M(\mathrm{^{10}B})+m_n-M(\mathrm{^{11}B})]c^2$。
质子分离能
$S_p=[M(\mathrm{^{10}Be})+m_H-M(\mathrm{^{11}B})]c^2$
（采用中性原子质量）。
两式相减
\[
 S_n-S_p=
 [M(\mathrm{^{10}B})+m_n
 -M(\mathrm{^{10}Be})-m_H]c^2.
\]
代原书质量超额：
$\Delta_{\mathrm{^{10}B}}=0.01294$、
$\Delta_{\mathrm{^{10}Be}}=0.01354$、
$\Delta_n=0.00867$、
$\Delta_H=0.00783$，
\[
 \boxed{S_n-S_p=
 (0.01294+0.00867-0.01354-0.00783)
 (931.5\,\mathrm{MeV})
 \simeq0.22\,\mathrm{MeV}}.
\]
\excheck 差异来自质子受库仑排斥的附加效应，
同时受娘核结构差异影响，不能仅归因于核子静质量。

\nuclearreactionproblem{258}{261}
$\mathrm{^{20}Ne}$、
$\alpha$ 与 $\mathrm{^{12}C}$
平均每核子结合能分别为
$8.03$、$7.07$、$7.68\,\mathrm{MeV}$。
求将$\mathrm{^{20}Ne}$分为
两个$\alpha$与一个$\mathrm{^{12}C}$所需能量。
\exsolution
初核的结合能
$E_i=20(8.03)\,\mathrm{MeV}$。
末态三个核的总结合能
$E_f=2\cdot4(7.07)+12(7.68)$ MeV。
分裂所需能量为失去的结合能差
\[
 \boxed{E_{\rm req}=E_i-E_f
 =20(8.03)-8(7.07)-12(7.68)
 \simeq11.9\,\mathrm{MeV}}.
\]
\excheck 结合能越大质量亏损越大；
若初态结合能大于产物和，
分裂必须吸收正能量。
% === END 39-nuclear-reactions-249-258.tex ===

% === BEGIN 39-nuclear-reactions-259-266.tex blob 8da592ec91c0e647bd45a8adc1f68e8b7b01e99c ===
% 原书第六篇6.6节259--266；
% 印刷261--262，PDF270--271，书末答案350（PDF359）。
\nuclearreactionproblem{259}{261}
用原子质量单位计算：
(a) $\mathrm{^{8}Li}$ 原子质量，其核结合能
$E_b=41.3\,\mathrm{MeV}$；
(b) $\mathrm{^{10}C}$ 原子核质量，其平均每核子
结合能为 $6.04\,\mathrm{MeV}$。
\exsolution
取 $1u\,c^2=931.5\,\mathrm{MeV}$，
原书氢原子质量 $m_H=(1+0.00783)u$、
中子质量 $m_n=(1+0.00867)u$。
(a) Li 有 $Z=3,A=8$，
\[
 \boxed{M_{\mathrm{^{8}Li\,atom}}=
 3m_H+5m_n-\frac{41.3\,\mathrm{MeV}}{c^2}
 \simeq8.0225u}.
\]
(b) C 有 $Z=6,A=10$，
结合能 $E_b=10(6.04)=60.4\,\mathrm{MeV}$。
求的是\emph{原子核}而不是原子，
所以应使用裸质子质量
$m_p=m_H-m_e$：
\[
 \boxed{M_{\mathrm{^{10}C\,nucleus}}=
 6(m_H-m_e)+4m_n-\frac{60.4\,\mathrm{MeV}}{c^2}
 \simeq10.0135u}.
\]
\excheck (b) 若直接代氢\emph{原子}质量而不扣去6个电子，
将错误地计算中性碳原子的质量。

\nuclearreactionproblem{260}{261}
某核反应
$A_1+A_2\to A_3+A_4$，
四核的结合能分别为 $E_1,E_2,E_3,E_4$。
求反应能 $Q$。
\exsolution
核反应能量守恒
$Q=(M_1+M_2-M_3-M_4)c^2$。
各核质量由
$M_ic^2=Z_im_pc^2+N_im_nc^2-E_i$。
在电荷数与核子数守恒下，
所有自由核子的总静能在反应前后抵消，
因此
\[
 \boxed{Q=(E_3+E_4)-(E_1+E_2)}.
\]
\excheck 产物的总结合能大于反应物时，
$Q>0$，说明放能；反之吸能。

\nuclearreactionproblem{261}{261--262}
假设每个 $\mathrm{^{235}U}$ 核裂变释放
$200\,\mathrm{MeV}$。
(a) 求1kg $\mathrm{^{235}U}$ 裂变释放的能量，
以及等能量、热值$30\,\mathrm{kJ/g}$的煤质量；
(b) 30千吨TNT当量的核爆炸对应多少裂变铀质量？
TNT热值为$4.1\,\mathrm{kJ/g}$。
\exsolution
(a) 1kg铀的核数
$N=(1000\,\mathrm g/235\,\mathrm{g/mol})N_A$，
释放总能量
\[
 \boxed{E=N(200\,\mathrm{MeV})
 \simeq8.21\times10^{13}\,\mathrm J
 =8.21\times10^{10}\,\mathrm{kJ}}.
\]
煤热值
$30\,\mathrm{kJ/g}=3.0\times10^7\,\mathrm{J/kg}$，
故
\[
 \boxed{m_{\rm coal}=
 \frac E{3.0\times10^7\,\mathrm{J/kg}}
 \simeq2.74\times10^6\,\mathrm{kg}}.
\]
(b) 一千吨TNT当量等于 $10^6\,\mathrm{kg}$ TNT，
30千吨即 $3.0\times10^7\,\mathrm{kg}$。
其能量
$E_{\rm TNT}=3.0\times10^7(4.1\times10^6)\,\mathrm J
=1.23\times10^{14}\,\mathrm J$。
因此等能铀质量
\[
 \boxed{m_U=\frac{E_{\rm TNT}}
 {8.21\times10^{13}\,\mathrm{J/kg}}
 \simeq1.50\,\mathrm{kg}}.
\]
\excheck 本版书末(a)煤质量写
$2.7\times10^8\,\mathrm{kg}$，
但依原题 $30\,\mathrm{kJ/g}$
计算应为$2.7\times10^6\,\mathrm{kg}$，
差100倍。两结果保留校勘记录。

\nuclearreactionproblem{262}{262}
由氘核聚变生成1.0g $\mathrm{^{4}He}$
时释放多少热？该热量相当于热值
$30\,\mathrm{kJ/g}$ 的多少煤？
\exsolution
反应
$2\,\mathrm{^{2}H}\to\mathrm{^{4}He}+\gamma$
（这里只比较核静质量）。
原书原子质量超额
$\Delta_D=0.01410u$、
$\Delta_\alpha=0.00260u$，
因此
\[
 Q=(2\Delta_D-\Delta_\alpha)c^2
 =0.02560u\,c^2
 \simeq23.85\,\mathrm{MeV}.
\]
1克氦核数
$N=(1.0\,\mathrm g/4\,\mathrm{g/mol})N_A$，
故
\[
 \boxed{E=NQ\simeq5.75\times10^{11}\,\mathrm J
 =5.75\times10^8\,\mathrm{kJ}}.
\]
煤热值
$30\,\mathrm{kJ/g}=3.0\times10^7\,\mathrm{J/kg}$，
\[
 \boxed{m_{\rm coal}\simeq1.92\times10^4\,\mathrm{kg}}.
\]
\excheck 原书书末写
$5.74\times10^7\,\mathrm{kJ}$，
但同一行给煤约$2\times10^4\,\mathrm{kg}$；
根据质量亏损及原题热值，能量应为
$5.74\times10^8\,\mathrm{kJ}$，书末相差10倍。

\nuclearreactionproblem{263}{262}
利用原子质量表，
求反应
$\mathrm{^{6}Li}+\mathrm{^{2}H}\to2\alpha$
每个核子释放的能量，
并与 $\mathrm{^{235}U}$ 每次裂变放200MeV时的
每核子放能量比较。
\exsolution
原书质量超额（单位$u$）
$\Delta_{\mathrm{^{6}Li}}=0.01513$、
$\Delta_D=0.01410$、
$\Delta_\alpha=0.00260$。
故
\[
 Q=[0.01513+0.01410-2(0.00260)]
 (931.5\,\mathrm{MeV})
 \simeq22.38\,\mathrm{MeV}.
\]
反应参与核子总数8，
所以每核子放能
\[
 \boxed{Q/8\simeq2.80\,\mathrm{MeV/nucleon}}.
\]
铀裂变每核子平均放能
$\boxed{200/235\simeq0.851\,\mathrm{MeV/nucleon}}$，
二者比约3.3。
\excheck 这里按\emph{参与反应的全部8个核子}
而不是按两个产物的\emph{个数}来归一化。

\nuclearreactionproblem{264}{262}
$\mathrm{^{7}Li}+p\to2\alpha$，
已知 $\mathrm{^{7}Li}$、
$\mathrm{^{4}He}$ 每核子结合能分别为
$5.60$ 与 $7.06\,\mathrm{MeV}$。
求核反应能。
\exsolution
反应初态总结合能
$E_i=7(5.60)+0$，自由质子结合能为零。
末态
$E_f=2\cdot4(7.06)$。
反应能是结合能差
\[
 \boxed{Q=E_f-E_i=
 8(7.06)-7(5.60)\simeq17.3\,\mathrm{MeV}}.
\]
\excheck 产物总结合能大于初态，
故该反应放能。

\nuclearreactionproblem{265}{262}
反应
$\mathrm{^{14}N}(\alpha,p)\mathrm{^{17}O}$
中入射 $\alpha$ 动能
$T_\alpha=4.0\,\mathrm{MeV}$，
出射质子动能
$T_p=2.09\,\mathrm{MeV}$，
质子与入射方向夹角
$\theta=60^\circ$。
求反应能 $Q$。
\exsolution
以实验室系静靶核为参考，
动量守恒
$\boldsymbol p_\alpha=\boldsymbol p_p+\boldsymbol p_O$，
所以
$p_O^2=p_\alpha^2+p_p^2-2p_\alpha p_p\cos\theta$。
非相对论动量
$p^2=2mT$，故反冲氧核动能
\[
 T_O=\frac{m_\alpha T_\alpha+m_p T_p
 -2\sqrt{m_\alpha m_pT_\alpha T_p}\cos\theta}{m_O}.
\]
近似质量数 $m_\alpha:m_p:m_O=4:1:17$，
得
$T_O\simeq0.724\,\mathrm{MeV}$。
根据动能能量守恒
$Q=T_p+T_O-T_\alpha$，
\[
 \boxed{Q\simeq2.09+0.724-4.00
 \simeq-1.19\,\mathrm{MeV}\approx-1.2\,\mathrm{MeV}}.
\]
\excheck 质子出射角相对\emph{入射束}，
故反冲项是
$-2p_\alpha p_p\cos\theta$；
该反应为吸能反应。

\nuclearreactionproblem{266}{262}
利用原子质量表计算：
(a)$\mathrm{^{7}Li}(p,n)\mathrm{^{7}Be}$；
(b)$\mathrm{^{9}Be}(n,\gamma)\mathrm{^{10}Be}$；
(c)$\mathrm{^{7}Li}(\alpha,n)\mathrm{^{10}B}$；
(d)$\mathrm{^{16}O}(d,\alpha)\mathrm{^{14}N}$
四个核反应的反应能。
\exsolution
对守恒总核子数的反应，
用\emph{中性原子质量超额}
求静质量差
$Q=(\sum\Delta_{\rm in}-\sum\Delta_{\rm out})c^2$，
其中
$1u\,c^2=931.5\,\mathrm{MeV}$。
从原书表取
$\Delta_H=0.00783$、
$\Delta_n=0.00867$、
$\Delta_{\mathrm{^{7}Li}}=0.01601$、
$\Delta_{\mathrm{^{7}Be}}=0.01693$、
$\Delta_{\mathrm{^{9}Be}}=0.01219$、
$\Delta_{\mathrm{^{10}Be}}=0.01354$、
$\Delta_{\mathrm{^{4}He}}=0.00260$、
$\Delta_{\mathrm{^{10}B}}=0.01294$、
$\Delta_{\mathrm{^{16}O}}=-0.00509$、
$\Delta_D=0.01410$、
$\Delta_{\mathrm{^{14}N}}=0.00307$，单位均为$u$。
于是
\[
 \boxed{Q_a=(0.01601+0.00783-0.00867-0.01693)c^2
 \simeq-1.64\,\mathrm{MeV}},
\]
\[
 \boxed{Q_b=(0.01219+0.00867-0.01354)c^2
 \simeq+6.82\,\mathrm{MeV}},
\]
\[
 \boxed{Q_c=(0.01601+0.00260-0.00867-0.01294)c^2
 \simeq-2.79\,\mathrm{MeV}},
\]
\[
 \boxed{Q_d=(-0.00509+0.01410-0.00260-0.00307)c^2
 \simeq+3.11\,\mathrm{MeV}}.
\]
\excheck $\gamma$ 静质量为零；
全部反应电荷守恒，
可以用同一中性原子质量超额记法互相抵消电子质量。
% === END 39-nuclear-reactions-259-266.tex ===

% === BEGIN 39-nuclear-reactions-267-280.tex blob cfb82c9d928e3765d6883547d1d1887d2c6369be ===
% Irodov, Part VI §6.6, original questions 267--280 (printed pp.262--263).
% Source: user-provided Arihant PDF pages 271--272; answers PDF page 359.
% Use SI units except energies explicitly measured in MeV.
\nuclearreactionproblem{267}{262}
极慢中子照射静止的$^{10}\mathrm B$核，
发生反应
$^{10}\mathrm B(n,\alpha){}^{7}\mathrm{Li}$。
利用原子质量表求两种产物分离时的速度。
\exsolution
用中性原子质量计算反应能，
\[
 Q=\bigl[M(^{10}\mathrm B)+m_n-
 M(^{7}\mathrm{Li})-M(^{4}\mathrm{He})\bigr]c^2
 \simeq2.79\,\mathrm{MeV}.
\]
入射中子初动量可忽略，末态两核动量等大反向。
由$p_\alpha=p_{\rm Li}$及
$Q=p^2(1/m_\alpha+1/m_{\rm Li})/2$，得
\[
 T_\alpha=Q\frac{m_{\rm Li}}{m_\alpha+m_{\rm Li}}
 \simeq\frac7{11}Q=1.78\,\mathrm{MeV},\quad
 T_{\rm Li}\simeq\frac4{11}Q=1.01\,\mathrm{MeV}.
\]
所以
\[
 \boxed{v_\alpha=\sqrt{2T_\alpha/m_\alpha}
 \simeq9.2\times10^6\,\mathrm{m/s}},\qquad
 \boxed{v_{\rm Li}\simeq5.3\times10^6\,\mathrm{m/s}}.
\]
\excheck 动量守恒要求$m_\alpha v_\alpha=m_{\rm Li}v_{\rm Li}$；
原书答案约$0.92\times10^7$与$0.53\times10^7\,\mathrm{m/s}$。

\nuclearreactionproblem{268}{262}
质子轰击静止$^{7}\mathrm{Li}$靶，
激发吸能反应$^{7}\mathrm{Li}(p,n){}^{7}\mathrm{Be}$。
在何种入射质子动能下，产物中子可恰好静止？
\exsolution
用原子质量求得反应能$Q\simeq-1.65\,\mathrm{MeV}$。
若末态中子静止，入射质子的全部动量
$p_p$由$^{7}\mathrm{Be}$承担。
初始动能$T=p_p^2/(2m_p)$，
反冲铍核动能$T_{\rm Be}=p_p^2/(2M_{\rm Be})
=(m_p/M_{\rm Be})T$。
能量守恒
$T+Q=T_{\rm Be}$，得
\[
 \boxed{T=\frac{-Q}{1-m_p/M_{\rm Be}}}
 \simeq\frac{1.65}{1-1/7}\,\mathrm{MeV}
 \simeq1.9\,\mathrm{MeV}.
\]
\excheck 此题条件是\emph{中子末态速度为零}，
并非一般吸能反应的最小阈值条件；
二者虽相近，物理含义不同。

\nuclearreactionproblem{269}{262}
动能$T_\alpha=5.3\,\mathrm{MeV}$的
$\alpha$粒子打静止$^{9}\mathrm{Be}$靶，
发生$^{9}\mathrm{Be}(\alpha,n){}^{12}\mathrm C$，
放能$Q=+5.7\,\mathrm{MeV}$。
求垂直入射束方向飞出的中子动能。
\exsolution
设末态中子动量垂直$\boldsymbol p_\alpha$，
反冲碳核动量为$\boldsymbol p_C=\boldsymbol p_\alpha-\boldsymbol p_n$。
故$p_C^2=p_\alpha^2+p_n^2$。
由非相对论关系$p_i^2=2m_iT_i$，
\[
 T_C=\frac{m_\alpha}{m_C}T_\alpha+
 \frac{m_n}{m_C}T_n.
\]
代入能量守恒$T_\alpha+Q=T_C+T_n$，
\[
 \boxed{T_n=
 \frac{Q+(1-m_\alpha/m_C)T_\alpha}
 {1+m_n/m_C}}
 \simeq\frac{5.7+(1-4/12)5.3}{1+1/12}
 =8.5\,\mathrm{MeV}.
\]
\excheck 碳核虽重仍有反冲动能，
不得把$T_n$直接设为$T_\alpha+Q=11.0$MeV。

\nuclearreactionproblem{270}{262}
动能$T_p=1.0\,\mathrm{MeV}$的质子
轰击静止$^{7}\mathrm{Li}$核，
发生$p+{}^{7}\mathrm{Li}\to2\alpha$。
如果两个$\alpha$粒子的运动方向
相对于入射束方向对称，
求各自动能及两速度夹角。
\exsolution
由质量亏损$Q\simeq17.3\,\mathrm{MeV}$。
在对称条件下两$\alpha$动能相同，
能量守恒给
\[
 \boxed{T_\alpha=(T_p+Q)/2\simeq9.1\,\mathrm{MeV}}.
\]
记两末态动量与入射方向的夹角均为$\phi$，
垂直分量互消，平行分量满足
$p_p=2p_\alpha\cos\phi$。
用$p_p^2=2m_pT_p$、
$p_\alpha^2=2m_\alpha T_\alpha$，
\[
 \cos\phi=
 \sqrt{\frac{m_pT_p}{4m_\alpha T_\alpha}},
 \qquad
 \boxed{\theta=2\phi\simeq170.5^\circ}.
\]
\excheck 接近背向的两条末态轨迹仍略向前倾斜，
合动量方向才能与入射质子一致。

\nuclearreactionproblem{271}{263}
质量$m$的非相对论入射粒子
轰击质量$M$的静止核，触发
反应能$Q<0$的吸能核反应。
证明其最小入射动能满足
$T_{\rm th}=(-Q)(1+m/M)$。
\exsolution
入射束实验室动能为$T$，初态相对运动的
质心系动能为
\[
 T_{\rm CM}=\frac{\mu}{m}T
 =\frac{M}{m+M}T,\qquad
 \mu=\frac{mM}{m+M}.
\]
在阈值处，反应末态各粒子在质心系均无相对动能，
故反应所需的最小质心动能恰为$-Q$。
因此
\[
 \frac{M}{m+M}T_{\rm th}=-Q,\qquad
 \boxed{T_{\rm th}=(-Q)\frac{m+M}{M}}.
\]
进一步令吸能量$\Delta=-Q>0$。
末态阈值时全部产物在质心系相对静止，
末态静质量总和是
$m+M+\Delta/c^2$。
若静止靶的静质量为$M$，
初态和末态的四动量不变量分别为
\[
 \begin{aligned}
 s_{\rm in}&=(m+M)^2c^4+2Mc^2T,\\
 s_{\rm th}&=\left[(m+M)c^2+\Delta\right]^2.
 \end{aligned}
\]
令$s_{\rm in}=s_{\rm th}$，得到
\[
 \boxed{T_{\rm th}^{\rm(rel)}
 =\Delta\left(1+\frac mM\right)
 +\frac{\Delta^2}{2Mc^2}}.
\]
所以经典结果是忽略
$\Delta^2/(2Mc^2)$后的首项。
\excheck 此为非相对论阈值公式，
若反应能相对靶核静能不可忽略，
应使用本题推导的四动量精确修正。
若$\Delta$接近重核静能，经典阈值
会系统性偏低；二体及多体末态都应把
最终所有静止质量计入同一不变量。

\nuclearreactionproblem{272}{263}
质子轰击静止氘核，
将其分裂为质子和中子。
氘核结合能$E_b=2.2\,\mathrm{MeV}$，
求入射质子的最小动能。
\exsolution
反应可写为$p+d\to p+p+n$，
总反应能$Q=-E_b$。
入射质子质量$m_p$、
静止靶氘核质量$m_d\simeq2m_p$。
由前题阈值公式
\[
 \boxed{T_{\rm th}\simeq
 E_b\left(1+\frac{m_p}{m_d}\right)
 \simeq2.2(1+\tfrac12)
 =3.3\,\mathrm{MeV}}.
\]
\excheck 不能只给$E_b=2.2$MeV，
因为实验室入射粒子需携带整个体系的质心平动能。

\nuclearreactionproblem{273}{263}
同一单能质子束分别轰击锂、铍靶。
观测到$^{7}\mathrm{Li}(p,n){}^{7}\mathrm{Be}$
（$Q=-1.65\,\mathrm{MeV}$）发生，
而$^{9}\mathrm{Be}(p,n){}^{9}\mathrm B$
（$Q=-1.85\,\mathrm{MeV}$）未发生。
求质子动能可能的范围。
\exsolution
阈值公式$T_{\rm th}=(-Q)(1+m_p/M)$：
\[
 T_{{\rm th},\mathrm{Li}}
 \simeq1.65\left(1+\frac17\right)
 =1.886\,\mathrm{MeV},
\]
\[
 T_{{\rm th},\mathrm{Be}}
 \simeq1.85\left(1+\frac19\right)
 =2.056\,\mathrm{MeV}.
\]
第一反应允许、第二反应不允许意味着
\[
 \boxed{1.89\,\mathrm{MeV}\lesssim T_p
 <2.06\,\mathrm{MeV}}.
\]
\excheck 实验上反应阈值附近产额可能很低；
这里“发生/不发生”按理想运动学允许性判断。

\nuclearreactionproblem{274}{263}
中子照射静止$^{11}\mathrm B$核，
激发$(n,\alpha)$反应的入射阈值为
$T_{\rm th}=4.0\,\mathrm{MeV}$。
求反应能$Q$。
\exsolution
由$T_{\rm th}=(-Q)(1+m_n/M_B)$，
$\;m_n\simeq u$、$M_B\simeq11u$，所以
\[
 \boxed{Q=-T_{\rm th}\frac{M_B}{M_B+m_n}
 \simeq-4.0\frac{11}{12}
 =-3.7\,\mathrm{MeV}}.
\]
\excheck 反应能为负，
符合必须以非零阈值粒子激发的事实。

\nuclearreactionproblem{275}{263}
计算静止$^{7}\mathrm{Li}$靶在
$(p,n)$、$(p,d)$两反应中所需的最小质子动能。
\exsolution
$(p,n)$通道为
$^{7}\mathrm{Li}+p\to{}^{7}\mathrm{Be}+n$，
质量亏损给$Q_n\simeq-1.65$MeV。
$(p,d)$通道为
$^{7}\mathrm{Li}+p\to{}^{6}\mathrm{Li}+d$，
原子质量给
$Q_d=[M(^{7}\mathrm{Li})+M(^1\mathrm H)
 -M(^{6}\mathrm{Li})-M(^2\mathrm H)]c^2
 \simeq-5.03$MeV。
两者都以锂为静止靶，
\[
 \boxed{T_{\mathrm{th},(p,n)}
 =-Q_n(1+1/7)\simeq1.88\,\mathrm{MeV}},
\]
\[
 \boxed{T_{\mathrm{th},(p,d)}
 =-Q_d(1+1/7)\simeq5.75\,\mathrm{MeV}}.
\]
\excheck 相同入射核与靶核使两个阈值具有相同的
运动学修正因子$8/7$。

\nuclearreactionproblem{276}{263}
使用原子质量表，
求引发$^{7}\mathrm{Li}(\alpha,n){}^{10}\mathrm B$
所需的最小$\alpha$粒子动能，
并求此时$^{10}\mathrm B$核的实验室速度。
\exsolution
质量亏损给反应能
$Q=[M(^{7}\mathrm{Li})+M(^{4}\mathrm{He})
-M(^{10}\mathrm B)-m_n]c^2\simeq-2.79$MeV。
在静止锂核前，$m_\alpha\simeq4u$、
$M_{\rm Li}\simeq7u$，因此
\[
 \boxed{T_{\rm th}=(-Q)(1+4/7)
 \simeq4.4\,\mathrm{MeV}}.
\]
阈值处所有产物在质心系相对静止，
因此$^{10}\mathrm B$的实验室速度
等于初始质心速度：
\[
 v_B=V_{\rm CM}
 =\frac{p_\alpha}{m_\alpha+M_{\rm Li}}
 =\frac{\sqrt{2m_\alpha T_{\rm th}}}{m_\alpha+M_{\rm Li}}
 \simeq\boxed{5.3\times10^6\,\mathrm{m/s}}.
\]
\excheck 书末印$5.3\times10^4$ m/s；
按本题$4.4$MeV能量与动量守恒必约为
$5.3\times10^6$ m/s，差100倍，须记入勘误。

\nuclearreactionproblem{277}{263}
动能$T_n=10\,\mathrm{MeV}$的中子
引发$^{12}\mathrm C(n,\alpha){}^{9}\mathrm{Be}$。
反应阈值$T_{\rm th}=6.17\,\mathrm{MeV}$，
求以$90^\circ$离开入射中子方向的
$\alpha$粒子动能。
\exsolution
由静止碳靶阈值
\[
 Q=-\frac{M_C}{m_n+M_C}T_{\rm th}
 \simeq-\frac{12}{13}(6.17)
 =-5.70\,\mathrm{MeV}.
\]
出射$\alpha$与入射中子动量垂直，
故反冲Be核动量模平方
$p_{\rm Be}^2=p_n^2+p_\alpha^2$。
因此
$T_{\rm Be}=(m_n/M_{\rm Be})T_n+
(m_\alpha/M_{\rm Be})T_\alpha$。
代入能量守恒$T_n+Q=T_\alpha+T_{\rm Be}$，
\[
 \boxed{T_\alpha=
 \frac{(1-m_n/M_{\rm Be})T_n+Q}
 {1+m_\alpha/M_{\rm Be}}
 \simeq\frac{(1-1/9)10-5.70}{1+4/9}
 \simeq2.2\,\mathrm{MeV}}.
\]
\excheck 若忘掉Be核反冲，
会把入射动能减阈值的差误作$\alpha$的动能。

\nuclearreactionproblem{278}{263}
在光致氘核解离
$\gamma+d\to n+p$中，
氘结合能$E_b=2.2\,\mathrm{MeV}$。
求光子反应阈值比结合能高百分之多少。
\exsolution
初态为静止氘核与入射光子，
末态质心需守恒入射光子动量。
严格不变量阈值关系
\[
 (M_dc^2+\varepsilon)^2-\varepsilon^2
 =(m_p+m_n)^2c^4
 \quad\Rightarrow\quad
 \varepsilon_{\rm th}
 =\frac{(m_p+m_n)^2-M_d^2}{2M_d}c^2.
\]
设$(m_p+m_n-M_d)c^2=E_b$，化简
\[
 \varepsilon_{\rm th}
 =E_b+\frac{E_b^2}{2M_dc^2}.
\]
故相对超过值
\[
 \boxed{\frac{\varepsilon_{\rm th}-E_b}{E_b}\times100\%
 =\frac{E_b}{2M_dc^2}\times100\%
 \simeq0.059\%\approx0.06\%}.
\]
\excheck 高出的一小部分供末态质心反冲；
光子动量不可忽略。

\nuclearreactionproblem{279}{263}
动能$T_p=1.5\,\mathrm{MeV}$的质子
被静止氘核俘获，形成$^{3}\mathrm{He}$。
求生成核的激发能。
\exsolution
氘核质量约为$2u$，
初态质心相对动能
$T_{\rm CM}=[m_d/(m_d+m_p)]T_p
\simeq(2/3)T_p=1.0$MeV。
质子俘获成$^{3}\mathrm{He}$的
反应能由结合能差给
$Q=B(^{3}\mathrm{He})-B(d)
\simeq7.72-2.22=5.50$MeV。
故复合核的激发能
\[
 \boxed{E^*=Q+T_{\rm CM}
 \simeq5.50+\frac23(1.50)
 =6.5\,\mathrm{MeV}}.
\]
\excheck 质心平动能不能转换为核内激发，
所以不应直接把实验室全部1.5MeV加在$Q$上。
更重要的是，该$^{3}$He复合核能量
$E^*\simeq6.5$MeV已经\emph{高于}
$p+d$分离阈值$Q\simeq5.5$MeV，
因此属于可向$p+d$等通道解离的连续谱
复合核状态，而不是一个稳定的束缚激发能级。
若讨论实际探测到的$^{3}$He生成事件，
还必须指定辐射俘获或其他退激产物以带走能量。

\nuclearreactionproblem{280}{263}
$^{13}\mathrm C(d,n){}^{14}\mathrm N$
的核反应产额在氘核实验室动能
$0.60,0.90,1.55,1.80\,\mathrm{MeV}$
时出现极大。利用原子质量表，
求中间复合核的四个激发能级。
\exsolution
该反应经复合$^{15}\mathrm N$核发生。
以$^{13}\mathrm C+d$相对于
$^{15}\mathrm N$基态的质量能差
\[
 S_d=[M(^{13}\mathrm C)+M(d)-M(^{15}\mathrm N)]c^2
 \simeq16.16\,\mathrm{MeV}
\]
表示氘的结合阈值。每一个共振入射能
在质心系贡献
$T_{\rm CM}=[M_C/(M_C+m_d)]T_d
\simeq13T_d/15$。
所以
\[
 \boxed{E_i^*=S_d+\frac{13}{15}T_i}
\]
逐一得到
\[
 \boxed{E_i^*\simeq
 16.7,\quad16.9,\quad17.5,\quad17.7\ \mathrm{MeV}}.
\]
\excheck 峰值是激发的复合核$^{15}$N能级，
不能把四个实验室入射动能直接当作核能级。
% === END 39-nuclear-reactions-267-280.tex ===

% === BEGIN 39-nuclear-reactions-281-290.tex blob edb1c07416c90baa00d07a6cbc9d5f7f124d4b1f ===
% Irodov Part VI §6.6, original 281--290 (section complete).
% Printed pp.263--265, supplied PDF pp.272--274; answers PDF pp.359--360.
% In the scanned source, problem 286 is misprinted as 296: numbering restored.
\nuclearreactionproblem{281}{263}
窄束热中子通过厚$d=0.50\,\mathrm{mm}$的镉板后，
强度衰减为原来的$1/\eta$，
$\eta=360$。求中子与镉核相互作用的
有效截面（用靶材料密度和原子量计算）。
\exsolution
对于被准直束移出的中子，
$J(d)=J_0e^{-n\sigma d}$，
其中$n=\rho N_A/M$。
因此
\[
 \boxed{\sigma=\frac{\ln\eta}{nd}
 =\frac{M\ln\eta}{\rho N_A d}}.
\]
取镉$\rho\simeq8.65\,\mathrm{g/cm^3}$、
$M\simeq112.4\,\mathrm{g/mol}$、
$d=0.050\,\mathrm{cm}$，
\[
 \boxed{\sigma\simeq2.54\times10^{-21}\,\mathrm{cm^2}
 =2.54\times10^3\,\mathrm b\approx2.5\,\mathrm{kb}}.
\]
\excheck $1\,\mathrm b=10^{-24}\,\mathrm{cm^2}$；
衰减率用自然对数，不能用常用对数。

\nuclearreactionproblem{282}{263}
窄束热中子经过厚$d=5.0\,\mathrm{cm}$的重水。
中子与氘、氧核的有效作用截面分别为
$\sigma_D=7.0\,\mathrm b$、
$\sigma_O=4.2\,\mathrm b$。
求透射束的衰减倍数。
\exsolution
每个$\mathrm D_2\mathrm O$分子含两个氘核、一个氧核。
设重水分子数密度
$n=\rho N_A/M_{\mathrm D_2O}$，
总宏观截面
$\Sigma=n(2\sigma_D+\sigma_O)$。
由指数衰减
\[
 \boxed{\frac{J_0}{J}=
 \exp[n(2\sigma_D+\sigma_O)d]}.
\]
取重水密度$\rho\simeq1.11\,\mathrm{g/cm^3}$、
摩尔质量$M_{\mathrm D_2O}\simeq20.0\,\mathrm{g/mol}$，
得到指数约$3.04$，
\[
 \boxed{J_0/J\simeq20.9\approx20}.
\]
\excheck 氘核有两个，故应使用$2\sigma_D$；
指数项的物质数密度是\emph{分子}浓度。

\nuclearreactionproblem{283}{264}
窄束热中子穿过厚$d=0.50\,\mathrm{cm}$的铁板，
吸收截面$\sigma_a=2.5\,\mathrm b$，
散射截面$\sigma_s=11\,\mathrm b$。
求因散射离开原准直束的中子份额。
\exsolution
设$\sigma_t=\sigma_a+\sigma_s$，
铁核数密度$n=\rho N_A/M$。
在深度$x$首次发生散射的概率密度为
$n\sigma_se^{-n\sigma_tx}\,dx$，
因此散射损失份额为
\[
 w_s=\int_0^d n\sigma_se^{-n\sigma_tx}dx
 =\boxed{\frac{\sigma_s}{\sigma_t}
 \left(1-e^{-n\sigma_t d}\right)}.
\]
取$\rho=7.87\,\mathrm{g/cm^3}$、
$M=55.85\,\mathrm{g/mol}$，
$n\sigma_td\simeq0.573$，
所以
$\boxed{w_s\simeq0.355\approx0.35}$。
\excheck 总束流损失与散射损失不相同；
首次相互作用可能是吸收而不是散射。

\nuclearreactionproblem{284}{264}
生产放射性核素的核反应产额可分别表述为
$w=N_{\rm react}/N_{\rm beam}$，
以及$k=A/N_{\rm beam}$，
即生成核素活度与入射粒子数之比。
(a) 由$w,k$求生成物半衰期；
(b) 质子束电流$I=10\,\mu\mathrm A$
照射锂靶$t=2.0\,\mathrm h$，
产生$^{7}\mathrm{Li}(p,n){}^{7}\mathrm{Be}$，
辐照完$^{7}\mathrm{Be}$活度
$A=1.35\times10^8\,\mathrm{Bq}$，
半衰期$T_{1/2}=53\,\mathrm d$，
求反应产额$w$。
\exsolution
(a) 定义给$A=\lambda N_{\rm react}$
（忽略辐照期间生成物的损失），
所以$k=A/N_{\rm beam}=\lambda w$。
故
\[
 \boxed{T_{1/2}=\frac{\ln2}{\lambda}
 =\frac{w\ln2}{k}}.
\]
(b) 入射质子总数$N_{\rm beam}=It/e$。
生成的放射性核数$N_{\rm react}=A/\lambda=AT_{1/2}/\ln2$，
从而
\[
 \boxed{w=
 \frac{AT_{1/2}e}{It\ln2}}.
\]
代入$T_{1/2}=53\times86400$秒、
$t=7200$秒和基本电荷，
得到
$\boxed{w\simeq1.99\times10^{-3}\approx2.0\times10^{-3}}$。
\excheck 对给定$53$天半衰期，
两小时内生成物的衰变不足以显著影响此近似。

\nuclearreactionproblem{285}{264}
通量密度$J=1.0\times10^{10}\,
\mathrm{cm^{-2}s^{-1}}$的热中子
垂直照射质量$m=10\,\mathrm{mg}$的
稳定$^{197}\mathrm{Au}$薄箔，
中子俘获截面$\sigma=98\,\mathrm b$，
生成半衰期$T_{1/2}=2.7\,\mathrm d$的
$\beta$活性$^{198}\mathrm{Au}$。
(a) 需要辐照多长时间，
才能使$^{197}\mathrm{Au}$核数减少1.0\%？
(b) 长时间辐照时$^{198}\mathrm{Au}$的极限核数？
\exsolution
令$\eta=0.010$。母核被中子俘获速率
$R=J\sigma N_{197}$，故
$N_{197}(t)=N_0e^{-J\sigma t}$。
于是
\[
 \boxed{t=\frac{-\ln(1-\eta)}{J\sigma}}.
\]
代$J\sigma=1.0\times10^{10}\cdot
98\times10^{-24}=9.8\times10^{-13}\,\mathrm{s^{-1}}$，
得
\[
 \boxed{t\simeq1.03\times10^{10}\,\mathrm s
 \simeq3.25\times10^{2}\,\mathrm{yr}}.
\]
(b) 在生成物建立接近放射性平衡的
几天时间内，稳定金核耗尽极微弱，
可令母核$N_{197}\simeq N_0=mN_A/M$。
生成物动力学
$dN_{198}/dt=J\sigma N_0-\lambda N_{198}$，
$\lambda=\ln2/T_{1/2}$。
平衡时
\[
 \boxed{N_{\max}\simeq\frac{J\sigma N_0}{\lambda}
 =\frac{J\sigma mN_A T_{1/2}}{M\ln2}
 \simeq1.01\times10^{13}}.
\]
为消除“饱和量”和严格最大值的差别，
再考虑靶核耗尽，置$a=J\sigma$。
严格的耦合微分方程是
\[
 \frac{dN_{197}}{dt}=-aN_{197},\qquad
 \frac{dN_{198}}{dt}=aN_{197}-\lambda N_{198},
\]
所以
\[
 N_{198}(t)=\frac{aN_0}{\lambda-a}
 \left(e^{-at}-e^{-\lambda t}\right).
\]
对时间求导并令其为零，得到
\[
 \boxed{t_*=\frac{\ln(\lambda/a)}{\lambda-a}
 \simeq58.1\,\mathrm d},\qquad
 \boxed{N_{198}(t_*)\simeq1.008\times10^{13}}.
\]
与上面忽略靶耗尽得到的饱和核数相比，
严格最大值仅低约$5\times10^{-6}$，故近似可靠。
\excheck 底本(a)书末印$3\times10^6$年；
按题面$J,\sigma$及1\%的下降量，
严格指数计算得约325年，
书末相差约$10^4$倍。
(b)的“最大数”是\emph{母核近似不耗尽时的
饱和量}，并非无限时间后真实母核耗尽的极限：
若任由母核完全转化，最终两种放射性核数均趋零。

\nuclearreactionproblem{286}{264}
薄箔初始每单位面积含$n$个稳定核，
以热中子通量密度$J$垂直辐照。
俘获截面$\sigma$，生成核素衰变常数$\lambda$。
推导单位面积放射性核数$N(t)$的积累规律。
\exsolution
薄箔中稳定靶核数近似不变且产物初始为零，
生成速率为$J\sigma n$，
衰变损失率为$\lambda N$，故
\[
 \frac{dN}{dt}=J\sigma n-\lambda N,\quad N(0)=0.
\]
积分因子$e^{\lambda t}$，
\[
 \boxed{N(t)=\frac{J\sigma n}{\lambda}
 (1-e^{-\lambda t})}.
\]
若不允许忽略靶核耗尽，
令$n_s(t)=n e^{-J\sigma t}$，
则更精细的表达式是
\[
 \boxed{N(t)=\frac{J\sigma n}{\lambda-J\sigma}
 \left(e^{-J\sigma t}-e^{-\lambda t}\right)}.
\]
\excheck \emph{原书扫描把本题题号误印为296}，
实际位于285与287之间，按连续编号应为286。
两个公式分别对应恒定靶量与显式耗尽模型。

\nuclearreactionproblem{287}{264}
一片$m=0.20\,\mathrm g$的金箔
受热中子照射$t=6.0\,\mathrm h$，
在照射结束后又过$\tau=12\,\mathrm h$，
测得金箔活度
$A=1.9\times10^7\,\mathrm{Bq}$。
生成放射性核的截面$\sigma=96\,\mathrm b$，
半衰期$T_{1/2}=2.7\,\mathrm d$。
求照射期间中子通量密度。
\exsolution
稳定$^{197}\mathrm{Au}$初始核数
$N_0=mN_A/197$。
辐照时产生的放射性核数
$N_{\rm end}=(J\sigma N_0/\lambda)
(1-e^{-\lambda t})$。
等待$\tau$后
$A=\lambda N_{\rm end}e^{-\lambda\tau}$。
所以
\[
 \boxed{J=\frac{A e^{\lambda\tau}}
 {\sigma N_0(1-e^{-\lambda t})}},\qquad
 \lambda=\frac{\ln2}{T_{1/2}}.
\]
代入$T_{1/2}=2.7\times86400$秒，
得到
\[
 \boxed{J\simeq5.9\times10^9\,
 \mathrm{cm^{-2}s^{-1}}\approx6\times10^9}.
\]
\excheck 辐照后等待时间使活度进一步下降，
求$J$时必须乘$e^{+\lambda\tau}$补偿此衰变。

\nuclearreactionproblem{288}{264}
初始有$N_0=1000$个中子，
每一代中子倍增系数$k=1.05$。
求第100代中子数。
\exsolution
将初始中子定义为第1代，
递推式$N_{j+1}=kN_j$，
故$N_j=N_0k^{j-1}$。
因此
\[
 \boxed{N_{100}=1000(1.05)^{99}
 \simeq1.25\times10^5\approx1.3\times10^5}.
\]
\excheck 指数是99而不是100，
因为初始的1000个中子已是第1代。

\nuclearreactionproblem{289}{264--265}
一台铀反应堆热功率
$P=100\,\mathrm{MW}$，
每次裂变平均产生$\nu=2.5$个中子，
释放能量$E=200\,\mathrm{MeV}$。
求单位时间产生的中子数。
\exsolution
核裂变每秒次数
$\dot N_{\rm fiss}=P/E$，
其中$200\,\mathrm{MeV}
=3.20435\times10^{-11}\,\mathrm J$。
故中子产生率
\[
 \boxed{\dot N_n=\nu\frac{P}{E}
 =2.5\frac{10^8\,\mathrm{J/s}}
 {3.20435\times10^{-11}\,\mathrm{J}}
 \simeq7.8\times10^{18}\,\mathrm{s^{-1}}}.
\]
\excheck 此数为所有裂变产生的中子总数，
不等于逸出堆芯的中子数或热中子数。

\nuclearreactionproblem{290}{265}
热中子反应堆中一代中子的平均寿命
原书印为$\tau=1.10\,\mathrm s$，
倍增系数$k=1.010$。
求(a)经过$t=1.0\,\mathrm{min}$时
中子数与功率增大倍数；
(b)使功率增加$e$倍的堆周期。
\exsolution
逐代倍增$N_j=N_0k^j$，
若一代平均时间为$\tau$，
则在近似连续化$j=t/\tau$时
\[
 \frac{N(t)}{N_0}=k^{t/\tau}
 =\exp\!\left(\frac{t\ln k}{\tau}\right).
\]
直接按扫描题面$\tau=1.10$秒、
$t=60$秒计算，
\[
 \boxed{\frac{N(60\,\mathrm s)}{N_0}
 =1.010^{60/1.10}\simeq1.72}.
\]
堆周期由$N(T)=eN_0$得
\[
 \boxed{T_{\rm reactor}=\frac{\tau}{\ln k}
 \simeq110.5\,\mathrm s}.
\]
若按书末近似写为$T\simeq\tau/(k-1)$，
仍为$\simeq110$秒。
\excheck 原书书末分别印增益约$2\times10^4$、
周期$10$秒，两者与扫描题面
$\tau=1.10$秒、$k=1.010$和$t=1.0$分钟
均不相容；本稿保留题面数据的物理一致解，
不私改任何单位或数字。
% === END 39-nuclear-reactions-281-290.tex ===
